\chapter{Urbanisation in India}

\section{Space, Place and the Inequality of Space}

\begin{KeyConcept}
\begin{itemize}
\item \textbf{Place} = the physical characteristics of a location (Delhi as the capital, with Mughal architecture) --- the concern of geographers and historians. \textbf{Space} = a location's \textbf{culture, politics and economy} --- the concern of sociologists.
\item There is an \textbf{inequality of space}: different places have different power. Delhi is more powerful because it is the seat of central politics and holds the powerful institutions --- a disruption in Delhi affects all of India, while a disruption in Bihar, the North-East or Kerala does not. (Irom Sharmila's protest in the North-East went unheard for years, but the farmers' protest in Delhi made the Prime Minister repeal the three laws.)
\item \textbf{Urbanisation is the most important manifestation of modernity}: urban centres are the centres of power, economy, politics and administration, and the \textbf{urban way of life is the aspiration} of young Indians (created through mass media and films) --- which is why agriculture/rural life is seen as unskilled and non-profitable. (This aspiration is propagated through \textbf{Gramsci's civil society} --- where the state's \textbf{hegemony} is created.)
\item Urban centres carry the dominant value of \textbf{individualism} --- which is why \textbf{violence (communal etc.) typically starts in urban centres} and then spreads to rural ones (where community life makes violence a last resort).
\end{itemize}
\end{KeyConcept}

\section{City and Country: Folk, Popular and High Culture}

\begin{KeyConcept}
\begin{itemize}
\item Before industrialisation, the \textbf{city} was the \textbf{marker of civilization} (the Indus Valley's advanced drainage and technology), while the \textbf{country} propagated \textbf{folk cultures} (homogeneous local cultures).
\item Three kinds of culture:
\item \textbf{Folk (low) culture} --- the oral, democratic culture of common people, learned through socialisation, with no fixed rules (village women singing at rituals). It spreads through the movement of women (marriage) and men (migration), producing \textbf{fusions}.
\item \textbf{High culture} --- rule-based, written, transmitted through the \textbf{guru-shishya parampara} (Hindustani classical music, classical dance, Bharat Muni's writings; art films; Margaret Atwood, Salman Rushdie).
\item \textbf{Popular culture} --- a middle path for \textbf{mass consumption}: it borrows from both folk and high culture, with no rules but wide reach (Bollywood, Chetan Bhagat, Punjabi songs simplified with Hindi/Urdu words). Urban centres are the centres of \textbf{popular culture}.
\end{itemize}
\end{KeyConcept}

\section{Redfield and Singer's Typology of Cities}

\begin{KeyConcept}
\begin{itemize}
\item \textbf{Robert Redfield and Milton Singer} studied Indian cities through \textbf{cultural processes} (not just economic/geographical), and classified cities into four types --- two pre-industrial, two post-industrial:
\item \textbf{Pre-industrial cities}: (1) \textbf{administrative-cultural cities} --- cities of the literati and indigenous bureaucracy (Varanasi, Allahabad, Kyoto, Lahore, Lhasa); (2) \textbf{cities of native commerce} --- weekly fairs and cattle markets providing goods to villages (Pushkar).
\item \textbf{Post-industrial cities}: (3) the \textbf{metropolis} --- cities of the worldwide managerial and entrepreneur class (London, New York; in India Bombay, Madras, Calcutta); (4) \textbf{cities of modern administration} --- built by rulers to gain legitimacy (New Delhi built near Mughal Delhi to erase Mughal symbolism; Washington DC).
\item \textbf{Metropolis vs cosmopolitan}: a metropolis has diverse cultures \textbf{living in clusters without necessarily interacting} (Bihari rickshaw-pullers, Oriya plumbers in Malviya Nagar); a \textbf{cosmopolitan} city has diverse cultures \textbf{interacting and creating a new, homogeneous culture} (New York). Many sociologists say \textbf{India has no truly cosmopolitan city} --- its cities remain divided by caste, class and religion.
\item The \textbf{role of the city}: in the pre-industrial phase, to \textbf{transform little traditions into great traditions} (folk music into classical music) on the basis of a \textbf{moral order}; in the post-industrial phase, cities run on a \textbf{technical order} (Bangalore = software, Dharavi = manufacturing, Delhi = services), are \textbf{uprooted from their indigenous culture}, and have a \textbf{central business district (CBD)} --- a marketplace of largely \textbf{impersonal relations} among diverse groups (e.g. Connaught Place, planned by Albert Mayer for Nehru's New Delhi).
\end{itemize}
\end{KeyConcept}

\begin{QuickRevision}
\begin{itemize}
\item Space (culture/politics/economy --- sociology) vs place (physical --- geography); inequality of space (Delhi > others; Sharmila vs farmers' protest).
\item Urbanisation = manifestation of modernity; urban way of life = aspiration; individualism $\rightarrow$ violence starts in cities.
\item Folk (oral, democratic), high (rule-based, guru-shishya), popular (mass consumption, Bollywood) culture.
\item Redfield \& Singer: 4 cities --- administrative-cultural (Varanasi, Kyoto), native commerce (Pushkar), metropolis (Bombay), modern administration (New Delhi); metropolis vs cosmopolitan; moral vs technical order; CBD (Connaught Place, impersonal relations).
\end{itemize}
\end{QuickRevision}

\section{Primary and Secondary Urbanisation}

\begin{KeyConcept}
\begin{itemize}
\item \textbf{Redfield and Singer's primary and secondary urbanisation}:
\item \textbf{Primary urbanisation} --- a pre-civilized \textbf{folk society} becomes a \textbf{peasant settlement} (after the invention of agriculture), develops a \textbf{sacred culture} (a new ritual life --- e.g. Malinowski's point that sea-going rituals changed), then \textbf{intellectuals} who interpret that sacred culture, then an \textbf{administration} --- producing the first (indigenous) urbanisation (Varanasi). No conflict, because everyone follows the same culture.
\item \textbf{Secondary urbanisation} --- when this society is \textbf{invaded by a powerful foreign power}, the local culture is dominated, and new rules are framed on the foreign culture. This produces \textbf{conflict as a permanent condition of cities}, and (through \textbf{syncretism}) a composite culture; it also creates \textbf{new identities} --- the ideas of the \textbf{minority, the marginalized, tolerance, enclaves}, and a new \textbf{intelligentsia} preaching a \textbf{cosmopolitan} identity (initially a negative term).
\end{itemize}
\end{KeyConcept}

\section{Consequences and Features of Urbanisation}

\begin{KeyConcept}
\begin{itemize}
\item The \textbf{consequences of urbanisation}: little traditions convert into great traditions; centres of administration and governance are established; \textbf{economic centres} are created (for indigenous products --- weekly markets --- and for foreign products); and cities become centres of \textbf{trade}.
\item The \textbf{features of urbanisation}: a \textbf{non-agricultural economy}; \textbf{trading and services} (including pilgrimage); \textbf{transport and communication networks}; \textbf{individualism} (you must trade with strangers of other cultures); \textbf{increased population density}; a \textbf{division of labour} (skilled craftsmen) and the idea of \textbf{merit}.
\item \textbf{Kingsley Davis}: urbanisation is an \textbf{index of the transformation from a rural/agricultural economy to a modern/industrial economy} --- a rise in population and density, and above all a \textbf{non-farm economy}.
\item \textbf{Slums}: initially workers lived with their owners (attached labour), but when industrial cities made the owner--worker relation \textbf{impersonal}, workers lived apart and the \textbf{slum culture} developed. (Sociologically, slums are the 'ocean of cultures' --- centres of political socialisation and hope --- to be discussed further.)
\end{itemize}
\end{KeyConcept}

\begin{QuickRevision}
\begin{itemize}
\item Primary urbanisation: folk society $\rightarrow$ peasant settlement $\rightarrow$ sacred culture $\rightarrow$ intellectuals $\rightarrow$ administration (Varanasi; no conflict).
\item Secondary urbanisation: foreign invasion $\rightarrow$ domination, conflict, syncretism, minority/marginalized/tolerance, cosmopolitan intelligentsia.
\item Consequences: little$\rightarrow$great tradition; administration; economic centres; trade.
\item Features: non-agricultural economy; services; transport; individualism; density; division of labour; merit.
\item Kingsley Davis: urbanisation = index of rural-to-industrial transformation (non-farm economy).
\item Slums: from attached labour to impersonal relations.
\end{itemize}
\end{QuickRevision}

